\documentclass[11pt,a4paper]{article}

% Essential packages
\usepackage[utf8]{inputenc}
\usepackage[T1]{fontenc}
\usepackage{geometry}
\usepackage{fancyhdr}
\usepackage{titlesec}
\usepackage{color}
\usepackage{listings}
\usepackage{graphicx}
\usepackage{float}
\usepackage{longtable}
\usepackage{array}
\usepackage{booktabs}
\usepackage{multirow}
\usepackage{adjustbox}
\usepackage{hyperref}
\usepackage{xcolor}
\usepackage{enumitem}
\usepackage{amsmath}
\usepackage{amsfonts}
\usepackage{amssymb}

% Page setup
\geometry{margin=2.5cm}
\pagestyle{fancy}
\fancyhf{}
\rhead{CryptoGL Architecture Guide}
\lhead{Project Improvement Recommendations}
\rfoot{Page \thepage}

% Title formatting
\titleformat{\section}{\Large\bfseries\color{blue}}{\thesection}{1em}{}
\titleformat{\subsection}{\large\bfseries\color{darkgray}}{\thesubsection}{1em}{}
\titleformat{\subsubsection}{\normalsize\bfseries\color{darkgray}}{\thesubsubsection}{1em}{}

% Colors
\definecolor{codegreen}{rgb}{0,0.6,0}
\definecolor{codegray}{rgb}{0.5,0.5,0.5}
\definecolor{codepurple}{rgb}{0.58,0,0.82}
\definecolor{backcolour}{rgb}{0.95,0.95,0.92}
\definecolor{problemred}{rgb}{0.8,0.2,0.2}
\definecolor{solutiongreen}{rgb}{0.2,0.6,0.2}
\definecolor{currentblue}{rgb}{0.2,0.2,0.8}

% Code listing setup
\lstset{
    backgroundcolor=\color{backcolour},   
    commentstyle=\color{codegreen},
    keywordstyle=\color{magenta},
    numberstyle=\tiny\color{codegray},
    stringstyle=\color{codepurple},
    basicstyle=\ttfamily\footnotesize,
    breakatwhitespace=false,         
    breaklines=true,                 
    captionpos=b,                    
    keepspaces=true,                 
    numbers=left,                    
    numbersep=5pt,                  
    showspaces=false,                
    showstringspaces=false,
    showtabs=false,                  
    tabsize=2,
    frame=single,
    language=C++,
    morekeywords={constexpr,noexcept,override,final,using,typename,template}
}

\begin{document}

\title{\Huge\textbf{CryptoGL Project Architecture\\Improvement Guide}}
\author{Architecture Optimization Recommendations}
\date{\today}

\maketitle

\begin{abstract}
This document provides a comprehensive analysis of the CryptoGL cryptographic library's current architecture and presents a prioritized roadmap for architectural improvements. The recommendations focus on modern C++17 practices, better code organization, improved maintainability, enhanced security, and optimized performance. Each recommendation includes current state analysis, problem identification, solution rationale, and implementation guidance.
\end{abstract}

\tableofcontents
\newpage

\section{Executive Summary}

The CryptoGL project is a comprehensive cryptographic library containing classical ciphers, symmetric/asymmetric ciphers, hash functions, and various cryptographic tools. While functionally complete, the current architecture has several areas for improvement in terms of modern C++ practices, code organization, maintainability, and security.

\subsection{Current Architecture Assessment}

\textbf{Strengths:}
\begin{itemize}
    \item Comprehensive coverage of cryptographic algorithms
    \item Good separation of concerns with header/implementation files
    \item Proper use of CMake build system
    \item Extensive unit test coverage
    \item Modular design with separate components
\end{itemize}

\textbf{Areas for Improvement:}
\begin{itemize}
    \item Flat file structure in \texttt{src/} directory
    \item Inconsistent naming conventions
    \item Mixed C++ and C-style code patterns
    \item Limited use of modern C++17 features
    \item No clear dependency management
    \item Missing abstraction layers
    \item Inconsistent error handling
\end{itemize}

\section{Priority 1: Foundation and Structure (Easy to Medium)}

\subsection{1.1 Directory Structure Reorganization}

\textbf{Current State:}
\begin{lstlisting}
src/
├── main.cpp
├── AES.cpp/hpp
├── DES.cpp/hpp
├── SHA2.cpp/hpp
├── Vector.hpp
├── String.hpp
├── Types.hpp
├── big_integers/
├── exceptions/
└── [100+ other files]
\end{lstlisting}

\textcolor{problemred}{\textbf{Problems with Current Structure:}}
\begin{itemize}
    \item All files in flat \texttt{src/} directory makes navigation difficult
    \item No logical grouping of related algorithms
    \item Difficult to find specific functionality
    \item Poor scalability for future additions
    \item No clear separation between core utilities and algorithms
\end{itemize}

\textcolor{solutiongreen}{\textbf{Why This Fix Works:}}
\begin{itemize}
    \item Logical grouping improves code discoverability
    \item Easier maintenance and refactoring
    \item Better dependency management
    \item Clearer project structure for new contributors
    \item Enables better build system organization
\end{itemize}

\textbf{Recommended Structure:}
\begin{lstlisting}
src/
├── core/
│   ├── types/
│   │   ├── Types.hpp
│   │   ├── Vector.hpp
│   │   └── String.hpp
│   ├── utils/
│   │   ├── Endian.hpp
│   │   ├── File.hpp
│   │   └── MathematicalTools.hpp
│   └── exceptions/
├── algorithms/
│   ├── block_ciphers/
│   │   ├── aes/
│   │   ├── des/
│   │   └── common/
│   ├── stream_ciphers/
│   ├── hash_functions/
│   ├── asymmetric/
│   └── classical/
├── modes/
│   └── BlockCipherModes.hpp
├── encoding/
│   └── Base64.hpp
└── main.cpp
\end{lstlisting}

\subsection{1.2 Namespace Organization}

\textbf{Current State:}
\begin{lstlisting}
// Mixed namespace usage
class AES { /* ... */ };
class Vector { /* ... */ };
namespace CryptoGL { /* ... */ };
\end{lstlisting}

\textcolor{problemred}{\textbf{Problems with Current Namespace Usage:}}
\begin{itemize}
    \item Inconsistent namespace usage across files
    \item Some classes in global namespace, others in CryptoGL namespace
    \item No logical grouping within namespaces
    \item Potential naming conflicts
\end{itemize}

\textcolor{solutiongreen}{\textbf{Why This Fix Works:}}
\begin{itemize}
    \item Prevents naming conflicts
    \item Provides logical grouping
    \item Improves code organization
    \item Makes dependencies explicit
\end{itemize}

\textbf{Recommended Namespace Structure:}
\begin{lstlisting}
namespace CryptoGL {
    namespace Core {
        namespace Types { /* Vector, String, etc. */ }
        namespace Utils { /* Endian, File, etc. */ }
        namespace Exceptions { /* All exception classes */ }
    }
    
    namespace Algorithms {
        namespace BlockCiphers { /* AES, DES, etc. */ }
        namespace StreamCiphers { /* RC4, Salsa20, etc. */ }
        namespace HashFunctions { /* SHA2, MD5, etc. */ }
        namespace Asymmetric { /* RSA, etc. */ }
        namespace Classical { /* Caesar, Vigenere, etc. */ }
    }
    
    namespace Modes { /* CBC, ECB, etc. */ }
    namespace Encoding { /* Base64, etc. */ }
}
\end{lstlisting}

\subsection{1.3 CMake Build System Modernization}

\textbf{Current State:}
\begin{lstlisting}
# Current CMakeLists.txt structure
set(SRC
   main.cpp
   AES.cpp
   DES.cpp
   # ... 100+ files listed manually
)

add_library(cryptoGL ${SRC})
\end{lstlisting}

\textcolor{problemred}{\textbf{Problems with Current Build System:}}
\begin{itemize}
    \item Manual file listing in CMakeLists.txt
    \item No modular component structure
    \item Difficult to build individual components
    \item No proper dependency management
    \item Hard to maintain as project grows
\end{itemize}

\textcolor{solutiongreen}{\textbf{Why This Fix Works:}}
\begin{itemize}
    \item Automatic file discovery reduces maintenance
    \item Modular structure enables selective building
    \item Better dependency tracking
    \item Easier to add new components
\end{itemize}

\textbf{Recommended CMake Structure:}
\begin{lstlisting}
# Root CMakeLists.txt
cmake_minimum_required(VERSION 3.15)
project(CryptoGL VERSION 1.0 LANGUAGES CXX)

set(CMAKE_CXX_STANDARD 17)
set(CMAKE_CXX_STANDARD_REQUIRED ON)

# Find all subdirectories automatically
file(GLOB_RECURSE SUBDIRS "src/*/CMakeLists.txt")
foreach(SUBDIR ${SUBDIRS})
    get_filename_component(SUBDIR_PATH ${SUBDIR} DIRECTORY)
    add_subdirectory(${SUBDIR_PATH})
endforeach()

# Main executable
add_executable(${PROJECT_NAME} src/main.cpp)
target_link_libraries(${PROJECT_NAME} 
    PRIVATE 
    CryptoGL::Core
    CryptoGL::Algorithms
    CryptoGL::Modes
)
\end{lstlisting}

\section{Priority 2: Modern C++17 Integration (Medium)}

\subsection{2.1 Smart Pointer Migration}

\textbf{Current State:}
\begin{lstlisting}
// Current raw pointer usage
class BlockCipher {
private:
    uint8_t* current_block;
    uint8_t* round_keys;
    
public:
    BlockCipher() {
        current_block = new uint8_t[16];
        round_keys = new uint8_t[176];
    }
    
    ~BlockCipher() {
        delete[] current_block;
        delete[] round_keys;
    }
};
\end{lstlisting}

\textcolor{problemred}{\textbf{Problems with Current Memory Management:}}
\begin{itemize}
    \item Manual memory management prone to leaks
    \item No exception safety
    \item Potential use-after-free bugs
    \item Inconsistent ownership semantics
\end{itemize}

\textcolor{solutiongreen}{\textbf{Why This Fix Works:}}
\begin{itemize}
    \item Automatic memory management
    \item Exception safety guaranteed
    \item Clear ownership semantics
    \item RAII compliance
\end{itemize}

\textbf{Recommended Implementation:}
\begin{lstlisting}
#include <memory>
#include <array>

class BlockCipher {
private:
    std::unique_ptr<std::array<uint8_t, 16>> current_block;
    std::unique_ptr<std::array<uint8_t, 176>> round_keys;
    
public:
    BlockCipher() 
        : current_block(std::make_unique<std::array<uint8_t, 16>>())
        , round_keys(std::make_unique<std::array<uint8_t, 176>>()) {
    }
    
    // No destructor needed - RAII handles cleanup
};
\end{lstlisting}

\subsection{2.2 Exception Handling Modernization}

\textbf{Current State:}
\begin{lstlisting}
// Current exception handling
class BadKeyLength : public std::exception {
public:
    const char* what() const noexcept override {
        return "Bad key length";
    }
};

// Usage
if (key.size() != 16) {
    throw BadKeyLength();
}
\end{lstlisting}

\textcolor{problemred}{\textbf{Problems with Current Exception Handling:}}
\begin{itemize}
    \item Basic exception classes with minimal information
    \item No error context or details
    \item Inconsistent error reporting
    \item No error categorization
\end{itemize}

\textcolor{solutiongreen}{\textbf{Why This Fix Works:}}
\begin{itemize}
    \item Rich error information for debugging
    \item Consistent error reporting across library
    \item Better error categorization
    \item Improved debugging experience
\end{itemize}

\textbf{Recommended Exception System:}
\begin{lstlisting}
#include <stdexcept>
#include <string>
#include <sstream>

namespace CryptoGL::Core::Exceptions {

enum class ErrorCode {
    INVALID_KEY_LENGTH,
    INVALID_BLOCK_SIZE,
    INVALID_ALGORITHM,
    INVALID_MODE,
    INTERNAL_ERROR
};

class CryptoException : public std::runtime_error {
private:
    ErrorCode code_;
    std::string context_;
    
public:
    CryptoException(ErrorCode code, const std::string& message, 
                   const std::string& context = "")
        : std::runtime_error(message), code_(code), context_(context) {}
    
    ErrorCode code() const noexcept { return code_; }
    const std::string& context() const noexcept { return context_; }
    
    std::string full_message() const {
        std::ostringstream oss;
        oss << "CryptoGL Error [" << static_cast<int>(code_) << "]: ";
        oss << what();
        if (!context_.empty()) {
            oss << " Context: " << context_;
        }
        return oss.str();
    }
};

class InvalidKeyLengthException : public CryptoException {
public:
    InvalidKeyLengthException(size_t expected, size_t actual, 
                            const std::string& algorithm = "")
        : CryptoException(ErrorCode::INVALID_KEY_LENGTH,
                         "Invalid key length",
                         "Expected: " + std::to_string(expected) + 
                         ", Got: " + std::to_string(actual) + 
                         (algorithm.empty() ? "" : " Algorithm: " + algorithm)) {}
};

} // namespace CryptoGL::Core::Exceptions
\end{lstlisting}

\subsection{2.3 Template Metaprogramming for Algorithm Selection}

\textbf{Current State:}
\begin{lstlisting}
// Current algorithm selection
enum class CipherType {
    AES,
    DES,
    BLOWFISH
};

class CipherFactory {
public:
    static std::unique_ptr<BlockCipher> create(CipherType type) {
        switch(type) {
            case CipherType::AES: return std::make_unique<AES>();
            case CipherType::DES: return std::make_unique<DES>();
            case CipherType::BLOWFISH: return std::make_unique<Blowfish>();
            default: throw std::invalid_argument("Unknown cipher type");
        }
    }
};
\end{lstlisting}

\textcolor{problemred}{\textbf{Problems with Current Approach:}}
\begin{itemize}
    \item Runtime overhead for algorithm selection
    \item No compile-time type safety
    \item Difficult to extend with new algorithms
    \item No compile-time optimization opportunities
\end{itemize}

\textcolor{solutiongreen}{\textbf{Why This Fix Works:}}
\begin{itemize}
    \item Compile-time algorithm selection
    \item Type safety guaranteed at compile time
    \item Better optimization opportunities
    \item Easier to extend with new algorithms
\end{itemize}

\textbf{Recommended Template-Based Approach:}
\begin{lstlisting}
#include <type_traits>

namespace CryptoGL::Algorithms {

// Algorithm traits
template<typename T>
struct CipherTraits;

template<>
struct CipherTraits<AES> {
    static constexpr size_t block_size = 16;
    static constexpr size_t key_size = 32;
    static constexpr const char* name = "AES";
};

template<>
struct CipherTraits<DES> {
    static constexpr size_t block_size = 8;
    static constexpr size_t key_size = 8;
    static constexpr const char* name = "DES";
};

// Compile-time algorithm selection
template<typename CipherType>
class CipherWrapper {
    static_assert(std::is_base_of_v<BlockCipher, CipherType>,
                  "CipherType must inherit from BlockCipher");
    
private:
    CipherType cipher_;
    
public:
    static constexpr size_t block_size = CipherTraits<CipherType>::block_size;
    static constexpr size_t key_size = CipherTraits<CipherType>::key_size;
    
    template<typename... Args>
    CipherWrapper(Args&&... args) 
        : cipher_(std::forward<Args>(args)...) {}
    
    // Delegate to underlying cipher
    template<typename... Args>
    auto encrypt(Args&&... args) -> decltype(cipher_.encrypt(std::forward<Args>(args)...)) {
        return cipher_.encrypt(std::forward<Args>(args)...);
    }
    
    template<typename... Args>
    auto decrypt(Args&&... args) -> decltype(cipher_.decrypt(std::forward<Args>(args)...)) {
        return cipher_.decrypt(std::forward<Args>(args)...);
    }
};

// Usage example
using AESWrapper = CipherWrapper<AES>;
using DESWrapper = CipherWrapper<DES>;

// Compile-time validation
static_assert(AESWrapper::block_size == 16, "AES block size must be 16");
static_assert(DESWrapper::block_size == 8, "DES block size must be 8");

} // namespace CryptoGL::Algorithms
\end{lstlisting}

\section{Priority 3: Advanced Architecture Patterns (Medium to Hard)}

\subsection{3.1 Policy-Based Design for Algorithm Configuration}

\textbf{Current State:}
\begin{lstlisting}
// Current hardcoded algorithm configuration
class AES {
private:
    static constexpr size_t rounds = 10;
    static constexpr size_t block_size = 16;
    
public:
    void encrypt(const Vector& plaintext, Vector& ciphertext) {
        // Fixed implementation
    }
};
\end{lstlisting}

\textcolor{problemred}{\textbf{Problems with Current Design:}}
\begin{itemize}
    \item Hardcoded algorithm parameters
    \item No flexibility for different configurations
    \item Difficult to test different variants
    \item No compile-time customization
\end{itemize}

\textcolor{solutiongreen}{\textbf{Why This Fix Works:}}
\begin{itemize}
    \item Flexible algorithm configuration
    \item Compile-time customization
    \item Better testability
    \item Reusable components
\end{itemize}

\textbf{Recommended Policy-Based Design:}
\begin{lstlisting}
#include <type_traits>

namespace CryptoGL::Policies {

// Key scheduling policy
template<typename T>
struct KeySchedulingPolicy;

template<>
struct KeySchedulingPolicy<AES> {
    static constexpr size_t rounds = 10;
    static constexpr size_t key_size = 32;
    
    template<typename KeyType>
    static auto schedule_key(const KeyType& key) {
        // AES-specific key scheduling
        return AES::expandKey(key);
    }
};

// Block processing policy
template<typename T>
struct BlockProcessingPolicy;

template<>
struct BlockProcessingPolicy<AES> {
    static constexpr size_t block_size = 16;
    
    template<typename BlockType>
    static void process_block(BlockType& block) {
        // AES-specific block processing
        AES::processBlock(block);
    }
};

// Main algorithm template
template<typename Algorithm, 
         typename KeyPolicy = KeySchedulingPolicy<Algorithm>,
         typename BlockPolicy = BlockProcessingPolicy<Algorithm>>
class ConfigurableCipher {
private:
    Algorithm algorithm_;
    
public:
    static constexpr size_t rounds = KeyPolicy::rounds;
    static constexpr size_t block_size = BlockPolicy::block_size;
    static constexpr size_t key_size = KeyPolicy::key_size;
    
    template<typename KeyType>
    void set_key(const KeyType& key) {
        auto scheduled_key = KeyPolicy::schedule_key(key);
        algorithm_.set_scheduled_key(scheduled_key);
    }
    
    template<typename BlockType>
    void encrypt_block(BlockType& block) {
        BlockPolicy::process_block(block);
    }
};

// Usage with different configurations
using StandardAES = ConfigurableCipher<AES>;
using CustomAES = ConfigurableCipher<AES, CustomKeyPolicy, CustomBlockPolicy>;

} // namespace CryptoGL::Policies
\end{lstlisting}

\subsection{3.2 Type-Safe Algorithm Interface}

\textbf{Current State:}
\begin{lstlisting}
// Current interface design
class BlockCipher {
public:
    virtual void encrypt(const Vector& plaintext, Vector& ciphertext) = 0;
    virtual void decrypt(const Vector& ciphertext, Vector& plaintext) = 0;
    virtual void setKey(const Vector& key) = 0;
};

class AES : public BlockCipher {
public:
    void encrypt(const Vector& plaintext, Vector& ciphertext) override;
    void decrypt(const Vector& ciphertext, Vector& plaintext) override;
    void setKey(const Vector& key) override;
};
\end{lstlisting}

\textcolor{problemred}{\textbf{Problems with Current Interface:}}
\begin{itemize}
    \item Runtime polymorphism overhead
    \item No compile-time type safety
    \item Difficult to optimize
    \item No compile-time validation of requirements
\end{itemize}

\textcolor{solutiongreen}{\textbf{Why This Fix Works:}}
\begin{itemize}
    \item Compile-time type safety
    \item Better optimization opportunities
    \item Clear interface contracts
    \item Easier to extend and maintain
\end{itemize}

\textbf{Recommended Type-Safe Interface:}
\begin{lstlisting}
#include <concepts>
#include <type_traits>

namespace CryptoGL::Concepts {

// Concept for block cipher requirements
template<typename T>
concept BlockCipher = requires(T t, const Vector& input, Vector& output, const Vector& key) {
    { T::block_size } -> std::convertible_to<size_t>;
    { T::key_size } -> std::convertible_to<size_t>;
    { t.encrypt(input, output) } -> std::same_as<void>;
    { t.decrypt(input, output) } -> std::same_as<void>;
    { t.set_key(key) } -> std::same_as<void>;
};

// Concept for stream cipher requirements
template<typename T>
concept StreamCipher = requires(T t, const Vector& input, Vector& output, const Vector& key) {
    { T::key_size } -> std::convertible_to<size_t>;
    { t.process(input, output) } -> std::same_as<void>;
    { t.set_key(key) } -> std::same_as<void>;
};

// Concept for hash function requirements
template<typename T>
concept HashFunction = requires(T t, const Vector& input, Vector& output) {
    { T::digest_size } -> std::convertible_to<size_t>;
    { t.hash(input, output) } -> std::same_as<void>;
    { t.reset() } -> std::same_as<void>;
};

} // namespace CryptoGL::Concepts

namespace CryptoGL::Algorithms {

// Type-safe algorithm interface
template<typename T>
requires CryptoGL::Concepts::BlockCipher<T>
class BlockCipherInterface {
private:
    T cipher_;
    
public:
    static constexpr size_t block_size = T::block_size;
    static constexpr size_t key_size = T::key_size;
    
    void encrypt(const Vector& plaintext, Vector& ciphertext) {
        static_assert(plaintext.size() % block_size == 0, 
                     "Plaintext size must be multiple of block size");
        cipher_.encrypt(plaintext, ciphertext);
    }
    
    void decrypt(const Vector& ciphertext, Vector& plaintext) {
        static_assert(ciphertext.size() % block_size == 0, 
                     "Ciphertext size must be multiple of block size");
        cipher_.decrypt(ciphertext, plaintext);
    }
    
    void set_key(const Vector& key) {
        static_assert(key.size() == key_size, 
                     "Key size must match algorithm requirements");
        cipher_.set_key(key);
    }
};

// Usage with compile-time validation
using AESInterface = BlockCipherInterface<AES>;
using DESInterface = BlockCipherInterface<DES>;

// This would cause a compile error if AES doesn't satisfy BlockCipher concept
static_assert(CryptoGL::Concepts::BlockCipher<AES>, "AES must satisfy BlockCipher concept");

} // namespace CryptoGL::Algorithms
\end{lstlisting}

\subsection{3.3 Dependency Injection and Configuration}

\textbf{Current State:}
\begin{lstlisting}
// Current hardcoded dependencies
class AES {
private:
    Vector current_block;
    Vector round_keys;
    std::array<uint8_t, 256> sbox;
    
public:
    AES() {
        // Hardcoded S-box initialization
        initialize_sbox();
    }
    
private:
    void initialize_sbox() {
        // Hardcoded S-box values
    }
};
\end{lstlisting}

\textcolor{problemred}{\textbf{Problems with Current Design:}}
\begin{itemize}
    \item Hardcoded dependencies
    \item Difficult to test with different configurations
    \item No flexibility for different implementations
    \item Tight coupling between components
\end{itemize}

\textcolor{solutiongreen}{\textbf{Why This Fix Works:}}
\begin{itemize}
    \item Flexible dependency injection
    \item Better testability
    \item Loose coupling
    \item Configurable behavior
\end{itemize}

\textbf{Recommended Dependency Injection Design:}
\begin{lstlisting}
#include <memory>
#include <functional>

namespace CryptoGL::Core {

// Configuration interface
class AlgorithmConfig {
public:
    virtual ~AlgorithmConfig() = default;
    virtual std::string name() const = 0;
    virtual size_t version() const = 0;
};

// S-box provider interface
class SBoxProvider {
public:
    virtual ~SBoxProvider() = default;
    virtual std::array<uint8_t, 256> get_sbox() const = 0;
    virtual std::array<uint8_t, 256> get_inverse_sbox() const = 0;
};

// Key scheduling interface
class KeyScheduler {
public:
    virtual ~KeyScheduler() = default;
    virtual Vector schedule_key(const Vector& key) const = 0;
};

// Default implementations
class DefaultAESConfig : public AlgorithmConfig {
public:
    std::string name() const override { return "AES"; }
    size_t version() const override { return 1; }
};

class DefaultAESSBox : public SBoxProvider {
public:
    std::array<uint8_t, 256> get_sbox() const override {
        // Return standard AES S-box
        return AES_SBOX;
    }
    
    std::array<uint8_t, 256> get_inverse_sbox() const override {
        // Return standard AES inverse S-box
        return AES_INVERSE_SBOX;
    }
};

class DefaultAESKeyScheduler : public KeyScheduler {
public:
    Vector schedule_key(const Vector& key) const override {
        // Standard AES key scheduling
        return AES::expandKey(key);
    }
};

// Configurable AES implementation
class ConfigurableAES {
private:
    std::shared_ptr<AlgorithmConfig> config_;
    std::shared_ptr<SBoxProvider> sbox_provider_;
    std::shared_ptr<KeyScheduler> key_scheduler_;
    
    Vector current_block;
    Vector round_keys;
    std::array<uint8_t, 256> sbox;
    std::array<uint8_t, 256> inverse_sbox;
    
public:
    ConfigurableAES(std::shared_ptr<AlgorithmConfig> config = nullptr,
                   std::shared_ptr<SBoxProvider> sbox_provider = nullptr,
                   std::shared_ptr<KeyScheduler> key_scheduler = nullptr)
        : config_(config ? config : std::make_shared<DefaultAESConfig>())
        , sbox_provider_(sbox_provider ? sbox_provider : std::make_shared<DefaultAESSBox>())
        , key_scheduler_(key_scheduler ? key_scheduler : std::make_shared<DefaultAESKeyScheduler>()) {
        
        initialize_sboxes();
    }
    
    void set_key(const Vector& key) {
        round_keys = key_scheduler_->schedule_key(key);
    }
    
    void encrypt(const Vector& plaintext, Vector& ciphertext) {
        // Use injected dependencies
        // Implementation using sbox_ and round_keys_
    }
    
private:
    void initialize_sboxes() {
        sbox = sbox_provider_->get_sbox();
        inverse_sbox = sbox_provider_->get_inverse_sbox();
    }
};

// Factory for creating configured algorithms
class AlgorithmFactory {
public:
    template<typename AlgorithmType>
    static std::unique_ptr<AlgorithmType> create(
        std::shared_ptr<AlgorithmConfig> config = nullptr,
        std::shared_ptr<SBoxProvider> sbox_provider = nullptr,
        std::shared_ptr<KeyScheduler> key_scheduler = nullptr) {
        
        return std::make_unique<AlgorithmType>(config, sbox_provider, key_scheduler);
    }
};

} // namespace CryptoGL::Core

// Usage examples
auto standard_aes = AlgorithmFactory::create<ConfigurableAES>();
auto custom_aes = AlgorithmFactory::create<ConfigurableAES>(
    std::make_shared<CustomAESConfig>(),
    std::make_shared<CustomSBoxProvider>(),
    std::make_shared<CustomKeyScheduler>()
);
\end{lstlisting}

\section{Priority 4: Performance and Security Architecture (Hard)}

\subsection{4.1 Memory Safety and Zero-Copy Operations}

\textbf{Current State:}
\begin{lstlisting}
// Current memory management
class Vector {
private:
    uint8_t* data_;
    size_t size_;
    
public:
    Vector(const Vector& other) {
        size_ = other.size_;
        data_ = new uint8_t[size_];
        std::copy(other.data_, other.data_ + size_, data_);
    }
    
    Vector& operator=(const Vector& other) {
        if (this != &other) {
            delete[] data_;
            size_ = other.size_;
            data_ = new uint8_t[size_];
            std::copy(other.data_, other.data_ + size_, data_);
        }
        return *this;
    }
};
\end{lstlisting}

\textcolor{problemred}{\textbf{Problems with Current Memory Management:}}
\begin{itemize}
    \item Unnecessary memory allocations and copies
    \item Potential memory leaks
    \item Performance overhead from copying
    \item No move semantics support
\end{itemize}

\textcolor{solutiongreen}{\textbf{Why This Fix Works:}}
\begin{itemize}
    \item Zero-copy operations where possible
    \item Efficient move semantics
    \item RAII-compliant memory management
    \item Better performance characteristics
\end{itemize}

\textbf{Recommended Memory-Safe Design:}
\begin{lstlisting}
#include <memory>
#include <span>
#include <string_view>

namespace CryptoGL::Core {

// Zero-copy view types
template<typename T>
using DataView = std::span<const T>;

template<typename T>
using MutableDataView = std::span<T>;

// Memory-safe vector with zero-copy operations
template<typename T>
class SecureVector {
private:
    std::unique_ptr<T[]> data_;
    size_t size_;
    size_t capacity_;
    
public:
    // Constructors
    SecureVector() : data_(nullptr), size_(0), capacity_(0) {}
    
    explicit SecureVector(size_t size) 
        : data_(std::make_unique<T[]>(size)), size_(size), capacity_(size) {}
    
    // Move constructor
    SecureVector(SecureVector&& other) noexcept
        : data_(std::move(other.data_))
        , size_(other.size_)
        , capacity_(other.capacity_) {
        other.size_ = 0;
        other.capacity_ = 0;
    }
    
    // Move assignment
    SecureVector& operator=(SecureVector&& other) noexcept {
        if (this != &other) {
            data_ = std::move(other.data_);
            size_ = other.size_;
            capacity_ = other.capacity_;
            other.size_ = 0;
            other.capacity_ = 0;
        }
        return *this;
    }
    
    // Zero-copy view operations
    DataView<T> view() const noexcept {
        return DataView<T>(data_.get(), size_);
    }
    
    MutableDataView<T> mutable_view() noexcept {
        return MutableDataView<T>(data_.get(), size_);
    }
    
    // Secure operations
    void secure_clear() noexcept {
        if (data_) {
            std::fill(data_.get(), data_.get() + size_, T{});
        }
        size_ = 0;
    }
    
    // RAII cleanup
    ~SecureVector() {
        secure_clear();
    }
    
    // Prevent copying for security
    SecureVector(const SecureVector&) = delete;
    SecureVector& operator=(const SecureVector&) = delete;
};

// Zero-copy algorithm interface
template<typename T>
concept ZeroCopyAlgorithm = requires(T t, DataView<uint8_t> input, MutableDataView<uint8_t> output) {
    { t.process(input, output) } -> std::same_as<void>;
    { t.block_size() } -> std::convertible_to<size_t>;
};

// Example zero-copy AES implementation
class ZeroCopyAES {
private:
    SecureVector<uint8_t> key_;
    std::array<uint8_t, 16> block_buffer_;
    
public:
    void set_key(DataView<uint8_t> key) {
        key_ = SecureVector<uint8_t>(key.size());
        std::copy(key.begin(), key.end(), key_.mutable_view().begin());
    }
    
    void encrypt_block(DataView<uint8_t> input, MutableDataView<uint8_t> output) {
        // Process block in-place without copying
        std::copy(input.begin(), input.end(), block_buffer_.begin());
        process_block(block_buffer_);
        std::copy(block_buffer_.begin(), block_buffer_.end(), output.begin());
    }
    
    size_t block_size() const { return 16; }
    
private:
    void process_block(std::array<uint8_t, 16>& block) {
        // AES block processing
    }
};

static_assert(ZeroCopyAlgorithm<ZeroCopyAES>, "ZeroCopyAES must satisfy ZeroCopyAlgorithm concept");

} // namespace CryptoGL::Core
\end{lstlisting}

\subsection{4.2 Constant-Time Security Architecture}

\textbf{Current State:}
\begin{lstlisting}
// Current timing-vulnerable code
bool String::operator==(const String& other) const {
    if (size() != other.size()) {
        return false;  // Early return creates timing side-channel
    }
    
    for (size_t i = 0; i < size(); ++i) {
        if (data_[i] != other.data_[i]) {
            return false;  // Early return creates timing side-channel
        }
    }
    return true;
}
\end{lstlisting}

\textcolor{problemred}{\textbf{Problems with Current Security:}}
\begin{itemize}
    \item Timing side-channels in comparisons
    \item Early returns reveal information
    \item Branch-dependent execution paths
    \item Vulnerable to timing attacks
\end{itemize}

\textcolor{solutiongreen}{\textbf{Why This Fix Works:}}
\begin{itemize}
    \item Constant-time operations
    \item No timing side-channels
    \item Secure against timing attacks
    \item Predictable execution time
\end{itemize}

\textbf{Recommended Constant-Time Design:}
\begin{lstlisting}
#include <array>
#include <bit>
#include <cstring>

namespace CryptoGL::Security {

// Constant-time utilities
class ConstantTime {
public:
    // Constant-time comparison
    template<typename T>
    static bool equal(const T* a, const T* b, size_t size) noexcept {
        volatile uint8_t result = 0;
        
        for (size_t i = 0; i < size; ++i) {
            result |= a[i] ^ b[i];
        }
        
        return result == 0;
    }
    
    // Constant-time selection
    template<typename T>
    static T select(bool condition, T a, T b) noexcept {
        // Use bit manipulation for constant-time selection
        T mask = static_cast<T>(-static_cast<int>(condition));
        return (a & mask) | (b & ~mask);
    }
    
    // Constant-time memory comparison
    static bool memcmp_constant_time(const void* a, const void* b, size_t size) noexcept {
        return equal(static_cast<const uint8_t*>(a), 
                    static_cast<const uint8_t*>(b), size);
    }
    
    // Constant-time string comparison
    static bool strcmp_constant_time(const std::string& a, const std::string& b) noexcept {
        size_t max_size = std::max(a.size(), b.size());
        
        // Pad shorter string with zeros for constant-time comparison
        std::array<uint8_t, 256> padded_a = {};
        std::array<uint8_t, 256> padded_b = {};
        
        std::copy(a.begin(), a.end(), padded_a.begin());
        std::copy(b.begin(), b.end(), padded_b.begin());
        
        return equal(padded_a.data(), padded_b.data(), max_size);
    }
};

// Secure string with constant-time operations
class SecureString {
private:
    std::unique_ptr<uint8_t[]> data_;
    size_t size_;
    
public:
    SecureString() : data_(nullptr), size_(0) {}
    
    explicit SecureString(const std::string& str) 
        : data_(std::make_unique<uint8_t[]>(str.size()))
        , size_(str.size()) {
        std::copy(str.begin(), str.end(), data_.get());
    }
    
    // Constant-time comparison
    bool constant_time_equals(const SecureString& other) const noexcept {
        if (size_ != other.size_) {
            // Still compare full buffers to avoid timing side-channel
            return ConstantTime::equal(data_.get(), other.data_.get(), 
                                     std::max(size_, other.size_));
        }
        return ConstantTime::equal(data_.get(), other.data_.get(), size_);
    }
    
    // Secure cleanup
    void secure_clear() noexcept {
        if (data_) {
            std::fill(data_.get(), data_.get() + size_, 0);
        }
    }
    
    ~SecureString() {
        secure_clear();
    }
    
    // Prevent copying for security
    SecureString(const SecureString&) = delete;
    SecureString& operator=(const SecureString&) = delete;
    
    // Allow moving
    SecureString(SecureString&&) = default;
    SecureString& operator=(SecureString&&) = default;
};

// Constant-time cryptographic operations
class ConstantTimeCrypto {
public:
    // Constant-time modular addition
    template<typename T>
    static T modular_add(T a, T b, T modulus) noexcept {
        T sum = a + b;
        T overflow = (sum < a) ? 1 : 0;
        T result = sum - (overflow * modulus);
        return result;
    }
    
    // Constant-time modular multiplication
    template<typename T>
    static T modular_multiply(T a, T b, T modulus) noexcept {
        T result = 0;
        T temp = a;
        
        for (size_t i = 0; i < sizeof(T) * 8; ++i) {
            T bit = (b >> i) & 1;
            result = ConstantTime::select(bit, 
                                        modular_add(result, temp, modulus), 
                                        result);
            temp = modular_add(temp, temp, modulus);
        }
        
        return result;
    }
    
    // Constant-time conditional swap
    template<typename T>
    static void conditional_swap(bool condition, T& a, T& b) noexcept {
        T mask = static_cast<T>(-static_cast<int>(condition));
        T temp = a ^ b;
        a ^= temp & mask;
        b ^= temp & mask;
    }
};

} // namespace CryptoGL::Security
\end{lstlisting}

\section{Implementation Roadmap and Priority Matrix}

\subsection{Architecture Improvement Priority Matrix}

\begin{table}[H]
\centering
\adjustbox{width=\textwidth}{
\begin{tabular}{|p{2cm}|p{3cm}|p{3cm}|p{3cm}|p{3cm}|}
\hline
\textbf{Priority Level} & \textbf{Effort Required} & \textbf{Impact} & \textbf{Implementation Time} & \textbf{Recommended Order} \\
\hline
Priority 1 & Low-Medium & High & 2-4 weeks & 1st \\
\hline
Priority 2 & Medium & High & 4-6 weeks & 2nd \\
\hline
Priority 3 & Medium-Hard & Medium-High & 6-8 weeks & 3rd \\
\hline
Priority 4 & Hard & High & 8-12 weeks & 4th \\
\hline
\end{tabular}
}
\caption{Architecture Improvement Priority Matrix}
\end{table}

\subsection{Detailed Implementation Timeline}

\begin{table}[H]
\centering
\adjustbox{width=\textwidth}{
\begin{tabular}{|p{2.5cm}|p{2cm}|p{2.5cm}|p{2.5cm}|p{2.5cm}|}
\hline
\textbf{Phase} & \textbf{Week} & \textbf{Primary Tasks} & \textbf{Secondary Tasks} & \textbf{Deliverables} \\
\hline
Foundation & 1-2 & Directory restructuring & Namespace organization & New project structure \\
\hline
Foundation & 3-4 & CMake modernization & Build system updates & Modernized build system \\
\hline
Modern C++ & 5-6 & Smart pointer migration & Exception handling & Memory-safe components \\
\hline
Modern C++ & 7-8 & Template metaprogramming & Concept definitions & Type-safe interfaces \\
\hline
Advanced Patterns & 9-10 & Policy-based design & Dependency injection & Configurable algorithms \\
\hline
Advanced Patterns & 11-12 & Type-safe interfaces & Algorithm factories & Extensible architecture \\
\hline
Security & 13-14 & Zero-copy operations & Memory safety & Performance optimizations \\
\hline
Security & 15-16 & Constant-time operations & Security hardening & Secure implementations \\
\hline
\end{tabular}
}
\caption{Detailed Implementation Timeline}
\end{table}

\subsection{Success Metrics and Validation}

\begin{table}[H]
\centering
\adjustbox{width=\textwidth}{
\begin{tabular}{|p{3cm}|p{3cm}|p{3cm}|p{3cm}|}
\hline
\textbf{Metric Category} & \textbf{Current State} & \textbf{Target State} & \textbf{Measurement Method} \\
\hline
Code Maintainability & Low (flat structure) & High (modular) & Cyclomatic complexity, file count per directory \\
\hline
Build Performance & Medium (manual files) & High (auto-discovery) & Build time, dependency resolution \\
\hline
Memory Safety & Medium (raw pointers) & High (smart pointers) & Static analysis, memory leak detection \\
\hline
Type Safety & Low (runtime checks) & High (compile-time) & Compile-time error count, concept validation \\
\hline
Security & Medium (timing attacks) & High (constant-time) & Timing analysis, security audit \\
\hline
Performance & Medium (copies) & High (zero-copy) & Benchmarking, profiling \\
\hline
\end{tabular}
}
\caption{Success Metrics and Validation Criteria}
\end{table}

\section{Conclusion}

This architecture improvement guide provides a comprehensive roadmap for modernizing the CryptoGL project. The recommendations are designed to:

\begin{itemize}
    \item Improve code maintainability and organization
    \item Enhance security through constant-time operations and memory safety
    \item Boost performance with zero-copy operations and modern C++17 features
    \item Increase type safety and compile-time validation
    \item Enable better testing and extensibility
\end{itemize}

The phased approach ensures that improvements can be implemented incrementally without disrupting existing functionality. Each phase builds upon the previous one, creating a solid foundation for future development.

\textbf{Key Benefits of Implementation:}
\begin{itemize}
    \item \textbf{Maintainability}: Modular structure and clear separation of concerns
    \item \textbf{Security}: Constant-time operations and memory-safe implementations
    \item \textbf{Performance}: Zero-copy operations and compile-time optimizations
    \item \textbf{Extensibility}: Policy-based design and dependency injection
    \item \textbf{Reliability}: Type-safe interfaces and comprehensive error handling
\end{itemize}

By following this roadmap, the CryptoGL project will evolve into a modern, secure, and high-performance cryptographic library that leverages the full power of C++17 and follows industry best practices.

\end{document} 